\chapter{Important Distributions}
In this appendix we will highlight the main probability density and mass functions that we will encounter most often during this course, as well as deriving and listing their useful properties. 

For now, here is a framework helped by AI to organize (as you can tell by the titles):

Each distribution we will write, plot, derive (some not all), moments, MLE, support, Moment generating function/characteristic function (some not all), cumulative (some not all), some special stuff for some (exponential period-frequency transformation, for example)

\section{Discrete distributions}
\subsubsection{Bernoulli: the foundation} %Sub sub sections give a nice bold but don't appear in the TOC... Win-win I'd say

\subsubsection{Binomial: For fixed trials}
Support: $k \in \{0, 1, \dots, n\}$ and $p \in [0, 1]$.

The probability of having exactly $k$ successes in $n$ independent Bernoulli trials is:

\begin{equation}
	\mathcal{B}_n(k;p) = \binom{n}{k}p^k\left(1-p\right)^{n-k}
\tag{Binomial pmf}
\label{pmf:binom}
\end{equation}

The Moment Generating Function (MGF) is the most efficient way to derive the raw moments:

	\[
	M_X(t) = \mathbb{E}[e^{tX}] = (1 - p + pe^t)^n
	\]
\todo{Not surem, and not sure how it brings us the moments...}
Which brings us to the mean and variance of $k$:

	\[
	\E\left[k\right] = np \quad \Var\left[k\right] = np(1-p)
	\]

The log-Likelihood for a sample of size $m$ (where each observation $k_i$ is the number of successes in $n$ trials) is:

	\[
	\ln L(p) = \sum_{i=1}^m \left[ \ln\binom{n}{k_i} + k_i\ln p + (n-k_i)\ln(1-p) \right]
	\]

which brings the MLE for $p$:

	\[
	\hat{p}_{MLE} = \frac{1}{mn} \sum_{i=1}^m k_i = \frac{\bar{k}}{n}
	\]

Which is an unbiased estimator with variance:

	\[
	\Var\left[\hat{p}_{MLE}\right] = \frac{1}{n^2}\Var[k] = \frac{np(1-p)}{n^2} = \frac{p(1-p)}{n}
	\]


Fisher Information: This measure tells us how much "information" $k$ carries about the unknown parameter $p$. For a single observation:

	\[
	\mathcal{I}(p) = -\E\left[ \frac{\partial^2 \ell}{\partial p^2} \right] = \frac{n}{p(1-p)}
	\] 

This shows us at a glance that $\hat{p}_{MLE}$ is the minimum variance estimator.

\subsubsection{Multinomial: The multivariate version of Binomial}
Add the example of the k-median here (distribution of the k-biggest value, it's an important exercise we did in class you'll see)
\subsubsection{Poisson: For occurrences over time/space}

\label{eq:poisson_mle}
We now want to find the maximum likelihood estimator (MLE) for a Poisson distribution. 
Given a sample of $N$ independent observations $\{x_i\}$, the total likelihood is:
\begin{align*}
	L_{\text{tot}}(\mu) &= \prod_{i=1}^N \frac{\mu^{x_i} e^{-\mu}}{x_i!}
\end{align*}
Taking the logarithm, we obtain the log-likelihood:
\begin{align*}
	\ln L_{\text{tot}}(\mu) 
	&= \sum_{i=1}^N \ln \left( \frac{\mu^{x_i} e^{-\mu}}{x_i!} \right) \\
	&= \sum_{i=1}^N x_i \ln \mu - \sum_{i=1}^N \mu - \sum_{i=1}^N \ln(x_i!)
\end{align*}
Since the last term does not depend on $\mu$, it can be neglected. Therefore,
\[
\ln L_{\text{tot}}(\mu) = \sum_{i=1}^N x_i \ln \mu - N\mu
\]

To find the MLE, we solve the \ref{eq:likelihood}:
\[
\frac{\partial \ln L_{\text{tot}}(\mu)}{\partial \mu} 
= \sum_{i=1}^N \frac{x_i}{\mu} - N = 0
\]
Evaluating at $\mu = \hat{\mu}$, we obtain:
\begin{equation*}
	\hat{\mu} = \frac{1}{N} \sum_{i=1}^N x_i
	\label{eq:poissonian_mle}
	\tag{Poisson MLE}
\end{equation*}

As expected, the MLE for the Poisson mean coincides with the sample mean. In the simple case of a single observation of $n$ events, we trivially get: $\hat\mu=n$.
\todo{I wrote this whole section using $\mu$ as the parameter instead of $\theta$ because I feel uneasy using $\theta$ for the expected events of a Poissonian, it just feels wrong. If we want to keep our convention going I can change it, no problem.}
%Negative Binomial / Geometric: For trials until success. Add if we ever use

\section{Continuos distributions}
\subsubsection{Uniform: The baseline for probability}

The general Uniform distribution $\mathcal{U}[l, u]$ between $l$ and $u$ with $l, u \in \mathbb{R}$, has \emph{pdf}:
\[
	p(x; l, u)=\begin{cases}
		\frac{1}{u-l} \; if \; x \in [l, v] \\
		0 \; \text{otherwise}
	\end{cases}
\]
One could write the same \emph{pdf} in terms of length of the interval $\delta$ and where it's centered $\mu$:
\[
p(x; \delta, \mu)=\begin{cases}
	\frac{1}{\delta} \; if \; x \in [\mu-\delta/2, \mu+\delta/2] \\
	0 \; \text{otherwise}
\end{cases}
\]

To make the following math easier we consider a Uniform $\mathcal{U}[0, m]$ (we can always reconduct ourself to this simple case trough a trivial change of variables).

We can evaluate the expected value of a uniform:
\begin{align*}
	\E[x]&=\int_{-\infty}^{+\infty}xp(x;0,m)\;dx=\int_{0}^{m}x\frac{1}{m}\;dx=\frac{m}{2}
\end{align*}
and for the variance:
\begin{align*}
	\operatorname{Var}[x]&=\E[x^2]-\E[x]^2=\int_{-\infty}^{+\infty}x^2p(x;0,m)\; dx - \frac{m^2}{4}\\
	&=\int_{0}^{m}x^2\frac{1}{m}\;dx- \frac{m^2}{4}=\frac{m^2}{3}-\frac{m^2}{4}=\frac{m^2}{12}
\end{align*}
For the characteristic function we get:
\[
	\phi(t)=\E[e^{itx}]=\frac{e^{imt}-1}{imt}
\]

We now want to compute the likelihood $L_x(m)$:
\[
	L_x(m)=\frac{1}{m}\mathbb{I}_{m\geq x}
\]
(see \ref{ex:uniform_likelihood}).
and for the Fisher information $I(m)$:
\[
	I(m)=\E \left[\left(\frac{\partial \ln L_x(m)}{\partial m}\right)^2 \right]= \E \left[\frac{1}{m^2} \right]=\frac{1}{m^2}
\]

Given $N$ observation ${x_i}$ following a uniform distribution, it is often useful to know how $x_{max}=\operatorname{max}[x_i]$ and $x_{min}=\operatorname{min}[x_i]$ are distributed.

We shall start with $x_{max}$. We can find it's distribution in two different but equivalent manner:
\begin{itemize}
	\item cumulative
	\item multinomial
\end{itemize}

\subsubsection{Normal (Gaussian): The king of statistics}

\subsubsection{Exponential: For time/space between events}

%Gamma: The generalized version of Exponential.  Add if we ever use

%Beta: For probabilities and proportions.  Add if we ever use

%Cauchy: The "pathological" example (useful to show when moments like mean/variance don't exist!).  Add if we ever use

\section{Sampling distributions}
%Student’s t: For small sample sizes.  Add if we ever use

\subsubsection{Chi-squared: For variances and goodness-of-fit}

%F-distribution: For comparing variances (ANOVA).  Add if we ever use